\clearpage\section{The appointed texts}

The Latin follows the 1962 Missal, pp.~398--400, nos.~1602--1611, with stress marks removed and ligatures expanded. The three English prayers are the historical human translation published by Cummiskey in 1861, pp.~431--433. Biblical English is the American 1899 Douay--Rheims canonical study text; it does not translate the Missal’s adaptations or its added formulas. Psalm numbering follows the Missal, with common modern equivalents in parentheses. These public-domain witnesses remain distinct from project-created exposition.

\subsection{Introit}\label{proper-introit}

\textit{Iustus es, Domine, et rectum iudicium tuum: fac cum servo tuo secundum misericordiam tuam.}

\textit{Ps. Beati immaculati in via: qui ambulant in lege Domini.}

\textit{V. Gloria Patri.}

\textbf{Douay--Rheims canonical study text.} \textsuperscript{137}SADE. Thou art just, O Lord: and thy judgment is right. \textsuperscript{124}Deal with thy servant according to thy mercy: and teach me thy justifications. \textsuperscript{1}ALEPH. Blessed are the undefiled in the way, who walk in the law of the Lord.

Ps.~118:137, 124; verse 1 (119). The Gloria Patri remains the Missal’s cue.

\subsection{Collect}\label{proper-collect}

\textit{Da, quaesumus, Domine, populo tuo diabolica vitare contagia: et te solum Deum pura mente sectari. Per Dominum nostrum.}

\textbf{Cummiskey 1861.} Grant, we beseech thee, O Lord, that thy people may avoid all contagion of the devil: and with a clean heart follow thee, the only true God. Thro’.

\subsection{Epistle}\label{proper-epistle}

\textit{Lectio Epistolae beati Pauli Apostoli ad Ephesios.}

\textit{Fratres: Obsecro vos ego vinctus in Domino, ut digne ambuletis vocatione, qua vocati estis, cum omni humilitate et mansuetudine, cum patientia, supportantes invicem in caritate, solliciti servare unitatem spiritus in vinculo pacis. Unum corpus, et unus spiritus, sicut vocati estis in una spe vocationis vestrae. Unus Dominus, una fides, unum baptisma. Unus Deus, et Pater omnium, qui est super omnes, et per omnia, et in omnibus nobis. Qui est benedictus in saecula saeculorum. Amen.}

\textbf{Douay--Rheims canonical study text.} \textsuperscript{1}I therefore, a prisoner in the Lord, beseech you that you walk worthy of the vocation in which you are called, \textsuperscript{2}With all humility and mildness, with patience, supporting one another in charity. \textsuperscript{3}Careful to keep the unity of the Spirit in the bond of peace. \textsuperscript{4}One body and one Spirit; as you are called in one hope of your calling. \textsuperscript{5}One Lord, one faith, one baptism. \textsuperscript{6}One God and Father of all, who is above all, and through all, and in us all.

Eph.~4:1--6. The Latin adds Fratres and the final doxology; these are absent from the canonical English above.

\subsection{Gradual}\label{proper-gradual}

\textit{Beata gens, cuius est Dominus Deus eorum: populus, quem elegit Dominus in hereditatem sibi.}

\textit{V. Verbo Domini caeli firmati sunt: et spiritu oris eius omnis virtus eorum.}

\textbf{Douay--Rheims canonical study text.} \textsuperscript{12}Blessed is the nation whose God is the Lord: the people whom he hath chosen for his inheritance. \textsuperscript{6}By the word of the Lord the heavens were established; and all the power of them by the spirit of his mouth:

Ps.~32:12, 6 (33), in the chant’s order.

\subsection{Alleluia}\label{proper-alleluia}

\textit{Alleluia, alleluia. V. Domine, exaudi orationem meam, et clamor meus ad te perveniat. Alleluia.}

\textbf{Douay--Rheims canonical study text.} \textsuperscript{2}Hear, O Lord, my prayer: and let my cry come to thee.

Ps.~101:2 (102:1). The biblical verse does not supply the Alleluia refrain.

\subsection{Gospel}\label{proper-gospel}

\textit{Sequentia sancti Evangelii secundum Matthaeum.}

\textit{In illo tempore: Accesserunt ad Iesum pharisaei: et interrogavit eum unus ex eis legis doctor, tentans eum: Magister, quod est mandatum magnum in lege? Ait illi Iesus: Diliges Dominum Deum tuum ex toto corde tuo, et in tota anima tua, et in tota mente tua. Hoc est maximum et primum mandatum. Secundum autem simile est huic: Diliges proximum tuum sicut teipsum. In his duobus mandatis universa lex pendet, et prophetae. Congregatis autem pharisaeis, interrogavit eos Iesus, dicens: Quid vobis videtur de Christo? cuius filius est? Dicunt ei: David. Ait illis: Quomodo ergo David in spiritu vocat eum Dominum, dicens: Dixit Dominus Domino meo, sede a dextris meis, donec ponam inimicos tuos scabellum pedum tuorum? Si ergo David vocat eum Dominum, quomodo filius eius est? Et nemo poterat ei respondere verbum: neque ausus fuit quisquam ex illa die eum amplius interrogare.}

\textit{Credo.}

\textbf{Douay--Rheims canonical study text.} \textsuperscript{34}But the Pharisees hearing that he had silenced the Sadducees, came together: \textsuperscript{35}And one of them, a doctor of the law, asking him, tempting him: \textsuperscript{36}Master, which is the greatest commandment in the law? \textsuperscript{37}Jesus said to him: Thou shalt love the Lord thy God with thy whole heart, and with thy whole soul, and with thy whole mind. \textsuperscript{38}This is the greatest and the first commandment. \textsuperscript{39}And the second is like to this: Thou shalt love thy neighbor as thyself. \textsuperscript{40}On these two commandments dependeth the whole law and the prophets. \textsuperscript{41}And the Pharisees being gathered together, Jesus asked them, \textsuperscript{42}Saying: What think you of Christ? whose son is he? They say to him: David’s. \textsuperscript{43}He saith to them: How then doth David in spirit call him Lord, saying: \textsuperscript{44}The Lord said to my Lord, Sit on my right hand, until I make thy enemies thy footstool? \textsuperscript{45}If David then call him Lord, how is he his son? \textsuperscript{46}And no man was able to answer him a word; neither durst any man from that day forth ask him any more questions.

Matt.~22:34--46. Canonical verse 34 supplies the Sadducee setting; the Missal replaces the opening. The appointed Gospel includes the full Son-of-David dialogue. Credo is the Missal’s following cue.

\subsection{Offertory}\label{proper-offertory}

\textit{Oravi Deum meum ego Daniel, dicens: Exaudi, Domine, preces servi tui: illumina faciem tuam super sanctuarium tuum: et propitius intende populum istum, super quem invocatum est nomen tuum, Deus.}

\textbf{Douay--Rheims canonical study text.} \textsuperscript{17}Now therefore, O our God, hear the supplication of thy servant, and his prayers: and shew thy face upon thy sanctuary which is desolate, for thy own sake. \textsuperscript{18}Incline, O my God, thy ear, and hear: open thy eyes, and see our desolation, and the city upon which thy name is called: for it is not for our justifications that we present our prayers before thy face, but for the multitude of thy tender mercies. \textsuperscript{19}O Lord, hear: O Lord, be appeased: hearken and do: delay not for thy own sake, O my God: because thy name is invocated upon thy city, and upon thy people.

Dan.~9:17--19 supplies context. The Latin chant introduces Daniel and rearranges selected petitions; the continuous English passage is not an English liturgical Offertory.

\subsection{Secret}\label{proper-secret}

\textit{Maiestatem tuam, Domine, suppliciter deprecamur: ut haec sancta, quae gerimus, et a praeteritis nos delictis exuant, et futuris. Per Dominum nostrum Iesum Christum, Filium tuum: Qui tecum vivit et regnat in unitate.}

\textit{Praefatio de Ssma Trinitate.}

\textbf{Cummiskey 1861.} We humbly beseech thy Majesty, O Lord, that the sacred mysteries we celebrate may cleanse us from all past offences, and from those we may hereafter be guilty of. Thro’.

The Trinity Preface is appointed; its full text is outside this proper-text reproduction. The historical English concerning future offences is explained in the commentary.

\subsection{Communion}\label{proper-communion}

\textit{Vovete et reddite Domino Deo vestro, omnes, qui in circuitu eius affertis munera: terribili, et ei qui aufert spiritum principum: terribili apud omnes reges terrae.}

\textbf{Douay--Rheims canonical study text.} \textsuperscript{12}Vow ye, and pay to the Lord your God: all you that are round about him bring presents. To him that is terrible, \textsuperscript{13}Even to him who taketh away the spirit of princes: to the terrible with the kings of the earth.

Ps.~75:12--13 (76:11--12).

\subsection{Postcommunion}\label{proper-postcommunion}

\textit{Sanctificationibus tuis, omnipotens Deus, et vitia nostra curentur, et remedia nobis aeterna proveniant. Per Dominum nostrum.}

\textbf{Cummiskey 1861.} May our vices be cured, O Almighty God! and an eternal remedy procured for us by these sacred mysteries. Thro’.
